\section{序論}
\label{sec:introduction}

正式標準は、どの製品同士が相互に機能するよう設計されるか、またどの企業が高コストの適応をせずに市場へ供給できるかを決める。しかし、政府が互いに非互換な標準を選んだとき、企業が受動的であり続ける必要はない。企業は製品をredesignし、追加的なtechnical specificationをサポートし、追加certificationを取得し、その他の方法によって複数の正式標準と自社製品を互換にすることができる。この私的能力は、harmonizationの経済学に基本的な緊張を生む。企業自身がformal incompatibilityを低コストで克服できるなら、public harmonizationは技術的には重要性が低下するように見える。しかし同じ私的応答は、exclusive standards arrangementを維持することで政府が得るrentも変え得る。したがって私的互換性（private compatibility）は、正式標準が作る技術的境界を弱めながら、その境界を維持する政治的誘因を変え得る。

formal compatibilityとeffective compatibilityの区別は経済的に重要である。標準協定には概念的に異なる2つの帰結があるからである。一方で制度上のmembershipを割り当て、他方で誰が同等の技術条件で競争できるかに影響する生産物市場の境界を作る。private workaroundは、第1の帰結を変えずに第2の帰結を緩和できる。例えば、追加的なtechnical specificationをサポートすれば、あるoutsiderは政府がstandards coalitionに加入していなくても、そのspecificationが適用される市場で互換企業と同じように競争できる場合がある。そのため正式coalitionは制度として存続しても、membershipを魅力的にしていたproduct-market rentの一部を失い得る。formal membershipとeffective compatibilityを同じ対象として扱うと、まさにこのfeedbackを排除してしまう。

本稿は、企業が標準の非互換性を私的に迂回できる能力が、正式な標準化coalitionのstabilityにどのような影響を与えるかを問う。本稿は、私的互換性が選択的侵食（selective erosion）のメカニズムを通じてexclusive regional standards unionを不安定化し得ることを示す。地域coalitionは当初、加盟国市場においてoutsiderをcompatibility groupから排除することで加盟国に利益を与え得るが、coalition加盟企業はoutsider marketへ販売するときには不利益を負ったままである。その後outsiderがbloc standardをサポートすることを有利と判断すると、その私的採用によって、加盟国市場でcoalition加盟国が持っていたcoalition-specific advantageは消える。しかしcoalition加盟企業はoutsider standardを採用していないため、outsider marketにおける相互的不利益は残る。formal membershipは変わらないが、formal coalitionに付随するcontinuation payoffは変わる。この非対称な変化は、地域標準化と国際標準化の間で加盟国政府の順位を逆転させ、その結果どのformal partitionがstableかを変え得る。

この議論を、formal institutional choiceとprivate technological adaptationを分離する3か国モデルで展開する。各国には1社の企業があり、企業は3つの分断された国内市場で数量競争を行う。政府は各国別標準を維持するか、3つの2か国地域標準化連合のいずれかを形成するか、共通の国際標準を採用するかを選べる。formal partitionは企業のbaseline compatibility relationと企業が生来的にサポートする標準を決める。そのpartitionが固定された後、企業は追加的な正式標準をサポートするために固定費用を支払うかを決める。1回の採用でその標準が適用されるすべての市場をカバーするので、私的互換性の収益はformal coalitionの地理的範囲に依存する。結果として得られるsupport patternを所与として、企業はcompatibility effectと、ローカル標準をサポートしない市場へ供給するためのadaptation costを伴うCournot product marketで競争する。その後政府は、国内consumer surplusと国内企業のworldwide profitの合計であるnational welfareを評価し、誘導されたcontinuation payoffがstrict coalition blockingを決める。

メカニズムは、地域coalition加盟国の国際標準化に対する相対利得を分解することで要約できる。private bypassがない場合、その差は
\begin{equation}
    \mathscr E_i-\mathscr D_i,
    \label{eq:introduction-selective-erosion}
\end{equation}
である。ここで $\mathscr E_i$ はcoalition加盟国市場で生じる加盟国の純排除余剰、$\mathscr D_i$ はoutsider marketにおける相互的不利益である。これらの対象にはproducer rentだけでなく、関連するnational-welfare consequenceが含まれる。outsiderがbloc standardを私的に採用すると、coalition marketのcompetitionはcomplete-compatibility configurationになる。したがってmember-market component $\mathscr E_i$ は消える。coalition加盟企業はoutsider standardを相互採用していないため、outsider-market component $\mathscr D_i$ は残る。したがって鍵となる変換は
\[
    \mathscr E_i-\mathscr D_i
    \quad\longrightarrow\quad
    -\mathscr D_i.
\]
である。$\mathscr E_i>\mathscr D_i>0$ なら、加盟国は迂回前には地域連合を好むが、迂回後には国際標準化を好む。私的適応はformal boundaryそのものを直接変えることなく、formal exclusionのpolitical valueを変える。

政治的効果の方向は技術的効果だけからは明らかではない。outsiderをregional member marketで互換にすると、regional regimeとinternational regimeのproduct-market differenceが縮小するので、より広いformal harmonizationの必要性が低下するように見える。しかしregional arrangementが加盟国にとって価値を持っていた理由の一部は、outsiderがそれらの市場から排除されていたことである。outsiderがその排除を迂回すると、coalition加盟国はそのadvantageを失う一方、なおincompatible outsider marketに直面する。したがって、あるformal boundaryの技術的な重要性を弱めるtechnologyが、そのboundaryを維持する魅力を低下させることがある。この意味で、technological substitutionはmultilateralizationに対する政府の相対的誘因の強化と両立する。

この分解はprivate compatibility costの役割も明確にする。固定費用は、あるformal regimeを機械的に有利にするexogenous welfare wedgeとしてメカニズムに入るのではない。それは、企業がどのprivate-adoption continuation stateを選ぶかを決める。regional standards unionでは、outsiderによるbloc standardへの投資は2つのmember marketで価値を持ち得る一方、coalition memberによるoutsider standardへの投資はoutsider market1つにしか作用しない。したがってprivate-adoption gameには、outsiderはbloc standardを採用するがどのmemberも相互採用しない非退化な範囲が存在する。その範囲ではprivate compatibilityの方向性が内生的である。同じformal regional partitionが、member marketではcomplete compatibilityを生みながら、outsider marketではreciprocal incompatibilityを残す。

headline resultは、このfirm-level responseをformal coalition stabilityへ写像する。対称的なMain Modelで、private circumventionが非活性のときにregional standardizationがcoalition memberにとって魅力的となるproduct-market parameterの非空open setを特定する。fixed compatibility costが高く $F>2T_A$ なら、regional unionでもseparate national standardsでも、どの企業も追加標準を私的にサポートしない。strict-blocking stable setはそのとき、3つすべての対称なregional standards unionからなる。この結果は一意のregional partnerを選ばない。対称性により、3つの2か国unionは同値だからである。

これに対し、固定費用が中間区間 $F_L<F<2T_A$ にあると、それぞれのregional outsiderだけがbloc standardを有利にサポートし、memberはoutsider standardをサポートせず、separate national standardsの下ではprivate adoptionは非活性のままである。選択的侵食により、各regional memberは現在のregional unionのcontinuation outcomeよりinternational standardizationを厳格に選好するようになる。outsiderもinternational standardizationを厳格に選好する。その結果grand coalitionがすべてのregional standards unionをblockし、international standardizationをblockするcoalitionは存在しない。したがって、本稿が維持するstrict-blocking conceptの下でinternational standardizationが一意にstableとなる。このtransitionはproduct-market parameterのopen setとcompatibility costの非退化なinterval上で生じ、threshold equalityによって駆動されない。さらに低いcompatibility costではより広いmultistandarding patternが生じ得るため、この比較の外でstable setが $F$ に対してmonotonicに変化するとは主張しない。

fixed-cost comparisonが有用なのは、formal institutional menuとunderlying product-market primitiveを変えずにprivate conductだけを変えるからである。したがって、private adaptationのfeasibilityからpublic cooperationのstabilityへ至るchannelを分離できる。adoption costが高いとき、企業は割り当てられたstandards structureを変えないのでformal exclusionは実質的なeconomic biteを持つ。中間費用では、private circumventionはregional boundaryの片側だけで有利となる。stable setの変化は、government preferenceの外生的変化や、国が直接international cooperationをより好むようになるという仮定ではなく、内生的なproduct-market continuation payoffによって駆動される。

副次的結果は、このstability mechanismをpartner selectionから切り離す。Main theoremはmarket sizeを等しくするため、private-compatibility-induced stability reversalにmarket-size heterogeneityは不要である。次に、その他のprimitiveを固定したまま、2国は同じ大きなmarket、3国目はより小さなmarketを持つ場合を考える。private adoptionがないことを条件とすると、大きい国はもう一方の大きい国をregional partnerとして選好する。小さいpartnerを大きいpartnerに置き換えると、market massがoutsider-profit blockからcoalition-member profit blockへ移るからである。この結果はpotential partnerを順位付けするものであり、coalition-stability theoremではなく、一意にstableなregional coalitionを意味しない。

本稿には互いに関連する3つの貢献がある。第1に、private compatibilityを政府のformal standards institutionに対する選好へ結び付けるmechanismとしてselective erosionを特定する。private compatibilityは、すべてのbilateral frictionを一様に減らすものとしてモデル化しない。adoptionはstrategicで、standard specificで、costlyである。bloc standardがoutsider standardとは異なる市場集合をカバーするため、均衡でのprivate responseには方向性が生じ得る。その方向性は政治的に重要である。outsiderは、coalition memberがbloc内で享受するexclusionary benefitを取り除きながら、bloc外でmemberが被る不利益を相互的に軽減しないからである。特定されたparameter regionでは、private compatibilityはformal harmonizationのmarket-access roleの一部を技術的に代替しながら、より広いharmonizationへの政府の相対的誘因を高め得る。後者はgovernment-incentive comparisonであり、一般的なcomplementarity theoremではない。

第2に、本モデルはfirm conductからinstitutional stabilityへのindustrial-organization channelを与える。formal standards regimeはbaseline compatibilityとmarginal adaptation costを決め、企業は追加標準をサポートすることによるproduct-market returnとfixed costを比較する。これらの意思決定は同じformal partitionの下でCournot competition、operating profit、consumer surplusを変える。政府はinstitutional labelそのものではなく、これらの内生的なcontinuation outcomeを評価する。因果連鎖は、企業のcompatibility choiceからproduct-market competition、coalition-specific rent、government incentive、formal coalition stabilityへ至る。このfeedbackこそが、本モデルをexogenous payoff tableを持つ単なるcoalition-formation exerciseではないものにしている。

第3に、この分析はregional-versus-multilateral standards questionにprivate-adaptation marginを加える。regional standards cooperationとmultilateral harmonizationは、競合するまたはsequentialなinstitutional arrangementとしてすでに研究されてきた。本稿の結果は新しいcoalition conceptや新しいfixed costを主張することには依存しない。むしろ同じformal regional partitionが、企業がそのtechnological boundaryを有利に迂回できるかどうかに応じて異なるpolitical continuation valueを持ち得る点にある。したがってcoalition gameは、post-formation firm behaviorが反応した後に評価される。この区別により、admissible government arrangementの集合が変わらなくても、private compatibility costの変化がformal coalition stabilityを変えることが可能になる。

これらの貢献は、複数の確立されたliteratureに基づく。standard、recognition、technical compatibilityをめぐるstrategic government choiceは、\citet{gandalShy2001}、\citet{barrettYang2001}、\citet{klimenko2009}、\citet{costinot2008} によって研究されている。regionalism questionに最も近い研究として、\citet{takaradaEtAl2020} はcoalition stabilityを用いた3か国Cournot modelでnational、regional、multilateral standardを比較し、複数標準の下で生産するfixed costも考える。\citet{kawabataTakarada2021} はregional standards harmonizationがmultilateral harmonizationを促進するか阻害するかを研究する。これらの論文は、strategic-policy environmentとregional standards coalitionの重要性の双方を確立している。本稿との違いは、formal partitionとgovernment continuation comparisonの間にstrategic private adoption stageを置くことである。

private-compatibility marginにも明確な先行研究がある。\citet{farrellSaloner1992} はconverterが本来非互換なtechnologyをbridgeできることを示し、\citet{doganogluWright2006} はnetwork marketでmultihomingがcompatibilityを代替し得ることを示す。\citet{fischerSerra2000} は異なるstandard levelで生産することに伴うfixed setup costを分析し、\citet{schmidtSteingress2022} はsunk adoption costの存在下で企業によるharmonized standardsのendogenous adoptionをモデル化する。したがって本稿の貢献は、企業が私的に標準差を克服できることや、その適応に費用がかかること自体ではない。追加標準supportの収益性がformal coalitionに依存し、その結果生じるprivate adoption patternがgovernment blocking incentiveの基礎となるcontinuation rentを変えることにある。

本稿はbroader economics of regulatory cooperationにも関連する。\citet{grossmanMcCalmanStaiger2021} は市場間でproduct tailoringを行うことにfixed costが伴うときのregulatory convergenceを研究する。\citet{parentiVannoorenberghe2025} はimperfect competitionとcountry-specific standardの下でregulatory cooperationとbloc formationを分析し、\citet{maggiMrazova2024} はregulatory diversityのfixed costがregulatory agreementなしでもspontaneous harmonizationを誘発し得ることを示す。これらの結果は、本稿の貢献を「harmonization plus fixed costs」と表現できない理由を補強する。本稿のより狭い対象は、formal regional coalitionが存続したままstrategic outsider bypassによって生じるselective continuation-payoff changeである。Section~\ref{sec:literature} でこれらの比較をさらに詳しく展開する。

stability concept自体は意図的にconventionalである。formal standards regimeをpartitionとして表し、coalitionが別のformal partitionを有利に誘導できるかをstrict blockingで評価する。これは \citet{hartKurz1983} に代表されるbroader coalition-formation traditionに基づく。実質的な新規性は、先行するindustrial-organization stageからcoalition gameへ供給されるpayoffにある。このことは結果の解釈にも節度を与える。stabilityはglobal efficiencyのstatementではない。intermediate-cost regionでinternational standardizationが一意にstableであることは、それがworld welfareを最大化することを示さず、regional unionのinstabilityもregional harmonizationがinefficientであることを意味しない。

本稿の構成は次のとおりである。Section~\ref{sec:literature} では、standards、compatibility、regulatory cooperation、coalitionに関する文献との関係で本メカニズムを位置付ける。Section~\ref{sec:model} ではmodelとtimingを提示する。Section~\ref{sec:product-market} でCournot building blockを導出し、Section~\ref{sec:private-compatibility} で関連するprivate-adoption incentiveを解く。Section~\ref{sec:selective-erosion} でselective-erosion decompositionを展開し、その後Section~\ref{sec:coalition-stability} でformal stability resultを確立する。Section~\ref{sec:secondary} では副次的なmarket-size partner-ranking resultを示す。Section~\ref{sec:discussion} でmechanismとscopeを議論し、Section~\ref{sec:conclusion} で結論を述べる。
